% ============================================================================
%  SPIE Section 4 -- RESULTS (tables, no figures)
%  Requires: \usepackage{booktabs, amsmath, amssymb}
%
%  DRAFTING STATUS -- READ BEFORE SUBMISSION
%  Every number below was measured with the policy reading the privileged
%  observation block (--privileged_actor sets privileged_obs_dim = 0, making
%  the actor's input slice a no-op). They are therefore TEACHER / oracle
%  numbers.
%    * With spie_methods.tex (teacher-student): Results must be REGENERATED
%      from the distilled student. Expect lower values.
%    * With spie_methods_asymmetric.tex: Results must come from the asymmetric
%      run on the second machine, whose policy is genuinely partially observed.
%  Do not ship either Methods variant paired with the numbers as they stand.
%  See SPIE_METHODS.md, inventory item R1.
% ============================================================================

\section{RESULTS}

Table~\ref{tab:clearance} reports the geometric audit of the generated anatomies. Reconstructing
the vessel surface from centerlines narrows the lumen by approximately 37\% of radius, which is
sufficient to make the $0.36$\,mm guidewire geometrically unable to pass at isolated stations:
before calibration, one of three sampled anatomies contained at least one station whose clearance
fell below the device radius. After calibration the generated anatomies reproduce the
patient-derived vessel almost exactly in median clearance ($2.14$\,mm versus $2.14$\,mm) while
remaining marginally tighter in the narrow tail (fifth percentile $1.13$\,mm versus $1.20$\,mm).
Generated anatomy is therefore no easier to traverse than the patient's own, and anatomical
variation is not obtained by widening the vessel.

\begin{table}[ht]
\caption{Lumen clearance along the navigated branch, measured to the nearest point on the
triangle surface. The calibrated generator matches the patient-derived reference in median
clearance while remaining marginally tighter at the fifth percentile.}
\label{tab:clearance}
\begin{center}
\begin{tabular}{lcccc}
\toprule
Surface & Median (mm) & 5th pct.\ (mm) & Min (mm) & Anatomies with a blocked station \\
\midrule
Patient-derived (reference) & 2.14 & 1.20 & 0.22 & 0 \\
Generated, uncalibrated     & 1.26 & 0.46 & 0.24 & 1 of 3 \\
Generated, calibrated       & 2.14 & 1.13 & 0.75 & 0 of 3 \\
\bottomrule
\end{tabular}
\end{center}
\end{table}

Table~\ref{tab:behavior} characterizes the two navigation behaviors during training for the two
training configurations, computed over all exploration episodes. The configurations differ only
in the potential-based shaping terms. The second configuration stalls more often (57.9\% versus
51.0\% of episodes) but resolves a substantially larger fraction of the stalls it encounters
(67.4\% versus 53.4\%), leaves far fewer stalls unresolved (32.6\% versus 46.6\%), and
essentially eliminates the catheter-advance pathology in which the catheter is driven ahead of
the guidewire (0.04 versus 0.34 of steps). Resolution by withdrawal of $1$--$8$\,mm---the
manoeuvre in which the device is eased back, reoriented, and re-advanced---rises from 22.6\% to
29.5\% of stall events.

\begin{table}[ht]
\caption{Behavioral characterization over all training-phase episodes. Stall resolution is
measured per stall event; episode success is measured per episode. The two configurations differ
only in the potential-based shaping terms.}
\label{tab:behavior}
\begin{center}
\begin{tabular}{lcc}
\toprule
Measure & Configuration A & Configuration B \\
\midrule
Exploration episodes analyzed          & 14{,}691 & 8{,}040 \\
Episodes containing a stall            & 51.0\%   & 57.9\%  \\
Stall events                           & 13{,}013 & 10{,}999 \\
Stall resolution rate (per event)      & 53.4\%   & \textbf{67.4\%} \\
\quad resolved, ${<}1$\,mm withdrawn   & 19.8\%   & 20.6\%  \\
\quad resolved, $1$--$8$\,mm withdrawn & 22.6\%   & \textbf{29.5\%} \\
\quad resolved, ${>}8$\,mm withdrawn   & 10.9\%   & 17.3\%  \\
\quad unresolved                       & 46.6\%   & \textbf{32.6\%} \\
Catheter-lead fraction of steps        & 0.34     & \textbf{0.04} \\
\bottomrule
\end{tabular}
\end{center}
\end{table}

Table~\ref{tab:synthetic} reports navigation success on 50 held-out synthetic anatomies, and
Table~\ref{tab:patient} on the patient-derived anatomy, both stratified by target depth. On
held-out synthetic anatomy the common carotid and mid-internal carotid segments are navigated
without failure (66 of 66 targets across both checkpoints); every remaining failure in the
evaluation is a cavernous-segment target. On the patient-derived anatomy, to which no model was
exposed during training, the best configuration reaches 75.5\% overall with the common carotid
solved completely and 90.2\% of mid-internal carotid targets reached.

\begin{table}[ht]
\caption{Navigation success on 50 held-out synthetic anatomies (98 episodes, deterministic
policy, 600-step limit, 5\,mm success threshold), stratified by target depth along the planned
path.}
\label{tab:synthetic}
\begin{center}
\begin{tabular}{lcccc}
\toprule
Model & Overall & Common carotid & Mid-ICA & Cavernous segment \\
\midrule
Configuration B, later checkpoint   & \textbf{84.7\%} (83/98) & 100\% (26/26) & 100\% (40/40) & 53.1\% (17/32) \\
Configuration B, earlier checkpoint & 83.7\% (82/98)          & 100\% (26/26) & 100\% (40/40) & 50.0\% (16/32) \\
\bottomrule
\end{tabular}
\end{center}
\end{table}

\begin{table}[ht]
\caption{Transfer to the patient-derived anatomy, evaluated on its original segmented surface
with all other settings unchanged. No model was exposed to this anatomy during training.}
\label{tab:patient}
\begin{center}
\begin{tabular}{lcccc}
\toprule
Model & Overall & Common carotid & Mid-ICA & Cavernous segment \\
\midrule
Configuration A, earlier checkpoint & 52.0\% (51/98)          & 92.6\% (25/27) & 58.5\% (24/41) & \phantom{0}6.7\% (2/30) \\
Configuration A, later checkpoint   & 63.3\% (62/98)          & 100\% (27/27)  & 65.9\% (27/41) & 26.7\% (8/30) \\
Configuration B, earlier checkpoint & 72.4\% (71/98)          & 88.9\% (24/27) & 63.4\% (26/41) & 70.0\% (21/30) \\
Configuration B, later checkpoint   & \textbf{75.5\%} (74/98) & 100\% (27/27)  & 90.2\% (37/41) & 33.3\% (10/30) \\
\bottomrule
\end{tabular}
\end{center}
\end{table}

Read together, Tables~\ref{tab:behavior}--\ref{tab:patient} indicate that the two training-phase
measures move together with transfer. Within Configuration B, the later checkpoint improves on
both measures simultaneously---episode success rises from 53.3\% to 59.2\% and stall resolution
from 44\% to 66\% over the corresponding training window---and improves on both evaluation sets,
from 83.7\% to 84.7\% on held-out synthetic anatomy and from 72.4\% to 75.5\% on the patient
anatomy. Across configurations the same ordering holds at comparable training-phase episode
success: Configuration B, whose stall-resolution rate was sustained through training, transfers
to the patient anatomy at 75.5\% against 63.3\% for Configuration A, whose resolution rate decayed
from 65\% to 41\% over the same run. Improvements confined to centerline following, without a
corresponding gain in stall resolution, produced the smaller transfer gain.

These comparisons should be read as consistent trends rather than as individually resolved
differences. At $n = 98$ the Wilson 95\% confidence intervals are
84.7\% $[76.3, 90.5]$, 83.7\% $[75.1, 89.7]$, 75.5\% $[66.1, 83.0]$ and 63.3\% $[53.4, 72.3]$: the
two checkpoints of Configuration B differ by a single episode, and the interval for the best
patient-anatomy result overlaps that of Configuration A. Episodes are additionally not fully
independent, and the initial device twist is not seeded, so repeat evaluations are not
bit-reproducible. The ordering is consistent across four checkpoints, two evaluation sets, and
three depth strata, but the individual margins are not established at this sample size.

Failures are confined almost entirely to the cavernous segment, where two near-$180^\circ$ bends
of a few millimetres' radius occur within a vessel of $2.5$--$3.4$\,mm diameter. The proximal and
mid-vessel course is solved on held-out synthetic anatomy, and the remaining problem is
depth-localized rather than uniformly distributed.
